% Generated draft: requires mathematical and source review.
\documentclass[11pt]{article}
\usepackage[margin=25mm]{geometry}
\usepackage{fontspec}
\setmainfont{Latin Modern Roman}
\usepackage{amsmath,amssymb,mathtools,unicode-math}
\setmathfont{Latin Modern Math}
\usepackage{enumitem,needspace,xurl,hyperref}
\setlength{\parindent}{0pt}
\setlength{\parskip}{5pt}
\setlength{\emergencystretch}{3em}
\newcommand{\N}{\mathbb{N}}
\newcommand{\Z}{\mathbb{Z}}
\newcommand{\Q}{\mathbb{Q}}
\newcommand{\R}{\mathbb{R}}
\newcommand{\C}{\mathbb{C}}
\newcommand{\F}{\mathbb{F}}
\newcommand{\A}{\mathbb{A}}
\newcommand{\PP}{\mathbb{P}}
\newcommand{\E}{\mathbb{E}}
\newcommand{\eps}{\varepsilon}
\newcommand{\dd}{\,\mathrm{d}}
\begin{document}
\section*{UnsolvedMath ID: 701031}
\textbf{Source problem code:} catalogue entry
\subsection{Definitions and mathematical statement}
Let \(V\) be a finite set, let \(k\geq 3\) be an integer, and let \(\mathcal L\subseteq\{L\subseteq V:|L|=k\}\). The members of \(V\) are called points and the members of \(\mathcal L\) are called lines. The pair \((V,\mathcal L)\) is a finite linear space in the sense of the supplied record if
\[
\forall x,y\in V,\quad x\neq y\ \Longrightarrow\ \bigl|\{L\in\mathcal L:x,y\in L\}\bigr|=1.
\]
Its automorphism group is
\[
\operatorname{Aut}(V,\mathcal L)=\{g\in\operatorname{Sym}(V):\{g(L):L\in\mathcal L\}=\mathcal L\},
\qquad g(L)=\{g(x):x\in L\}.
\]
A finite group \(G\) is almost simple if there is a finite nonabelian simple group \(T\) such that
\[
T\leq G\leq\operatorname{Aut}(T),
\]
where \(T\) is identified with its group of inner automorphisms. A subgroup \(G\leq\operatorname{Aut}(V,\mathcal L)\) is point-transitive if it acts transitively on \(V\), and line-transitive if it acts transitively on \(\mathcal L\).

\textbf{Classification problem.} Classify the finite linear spaces \((V,\mathcal L)\) admitting a subgroup \(G\leq\operatorname{Aut}(V,\mathcal L)\) that is almost simple, point-transitive, and line-transitive. The natural equivalence for this classification is incidence isomorphism: a bijection of point sets that maps the collection of lines onto the collection of lines. The source asks to classify the spaces admitting such a group, rather than explicitly asking to classify all choices of the group action.

\textbf{Boundary qualification.} The supplied definition does not explicitly require \(\mathcal L\neq\varnothing\), \(|V|\geq k\), or more than one line. No such additional restriction has been imposed here. The intended treatment of empty or otherwise degenerate incidence structures, including the convention for transitivity on an empty set, requires clarification.

\textbf{Supplied contextual qualifications.} For a point action of \(G\), point-quasiprimitivity means that every normal subgroup \(N\trianglelefteq G\) with \(N\neq 1\) is transitive on \(V\). The record attributes to A. Camina, P. Neumann, and C. E. Praeger a solution of the classification problem when \(T\) is an alternating group. It also reports that a line-transitive, point-quasiprimitive automorphism group of a finite linear space is almost simple or affine, citing A. Camina and C. E. Praeger (2001).

The remarks labelled 2005 report treatments of the simple groups \(\operatorname{PSU}(3,q)\), \(\operatorname{PSL}(2,q)\), \(\operatorname{Sz}(q)\), and \(\operatorname{Ree}(q)\), explicitly distinguishing those treatments from treatments of all almost simple groups with the corresponding simple socle. The same remarks report a treatment of sporadic simple groups. The record supplies no parameter ranges for these family-specific results.

The remark labelled 2009 reports that a point-quasiprimitive automorphism group of a non-Desarguesian projective plane would be affine; its cited publication is dated 2016. The remark labelled 2013 reports that the problem is solved for large-dimensional classical groups, citing a 2008 paper, without specifying a dimension threshold or the precise scope of that result. These are supplied historical reports, not independently verified conclusions or additional hypotheses of the classification problem. The registry labels the problem open; that label is supplied curation rather than an independently verified current literature assessment.
\subsection{Short English statement}
Classify finite linear spaces with a common line size \(k\geq 3\) that admit an almost simple automorphism group transitive on both points and lines. Preserve the distinction between results for a simple group itself and results for all almost simple groups with that socle. Degenerate cases and the exact scope of the reported partial results require review.
\subsection{Sources}
\begin{itemize}
\item Local UnsolvedMath registry, permanent ID \(701031\), exact source problem code \texttt{catalogue entry}, title \emph{Kourovka Notebook Problem 14.46}; supplied upstream revision \texttt{d371ee917ac9647fbb06cd6fc529deec6f3dda82}. The supplied record was created and updated on \texttt{2026-10-16T14:00:00Z} and carries the status label \texttt{open}.
\item \emph{The Kourovka Notebook (Unsolved Problems in Group Theory)}, Problem 14.46, proposed by A. Camina and C. E. Praeger in the 14th issue (1999). The supplied background identifies the 21st issue edition (2026), \texttt{arXiv:1401.0300v48}. The collection description names E. I. Khukhro and V. D. Mazurov as editors.
\item The supplied statement attributes the alternating-group case to A. Camina, P. Neumann, and C. E. Praeger; it provides no separate bibliographic citation for this attribution.
\item A. Camina and C. E. Praeger, \emph{Aequat. Math.} \textbf{61} (2001), 221--232. Supplied reference for the almost-simple-or-affine conclusion under line transitivity and point quasiprimitivity.
\item W. Liu, \emph{Linear Algebra Appl.} \textbf{374} (2003), 291--305. Supplied 2005 remark concerning \(\operatorname{PSU}(3,q)\).
\item W. Liu, \emph{J. Combin. Theory (A)} \textbf{103} (2003), 209--222. Supplied 2005 remark concerning \(\operatorname{PSL}(2,q)\).
\item W. Liu, \emph{Discrete Math.} \textbf{269} (2003), 181--190. Supplied 2005 remark concerning \(\operatorname{Sz}(q)\).
\item W. Liu, \emph{Europ. J. Combin.} \textbf{25} (2004), 311--325. Supplied 2005 remark concerning \(\operatorname{Ree}(q)\).
\item A. R. Camina and F. Spiezia, \emph{J. Combin. Des.} \textbf{8} (2000), 353--362. Supplied reference concerning sporadic simple groups.
\item N. Gill, \emph{Trans. Amer. Math. Soc.} \textbf{368} (2016), 3017--3057. Supplied reference in the remark labelled 2009 concerning non-Desarguesian projective planes.
\item A. Camina, N. Gill, and A. E. Zalesski, \emph{Bull. Belg. Math. Soc. Simon Stevin} \textbf{15}, no. 4 (2008), 705--731. Supplied reference in the remark labelled 2013 concerning large-dimensional classical groups.
\item All references and historical assessments above are reproduced from the supplied record; they have not been independently checked. No external source links or separate literature-verification date were supplied.
\end{itemize}
\end{document}
